% ──────────────────────────────────────────────────────────────────────
%  INFORME: Implementaci\'on de Autenticaci\'on Multifactor (MFA)
%  Curso: Seguridad Inform\'atica — Universidad de Pamplona
% ──────────────────────────────────────────────────────────────────────
\documentclass[12pt,letterpaper]{article}

\usepackage[T1]{fontenc}
\usepackage{berasans}
\renewcommand{\familydefault}{\sfdefault}
\usepackage{babel}
\babelprovide[import=es, main]{spanish}
\usepackage{geometry}
\geometry{margin=2.2cm}
\usepackage{setspace}
\singlespacing
\usepackage{graphicx}
\usepackage{tabularray}
\UseTblrLibrary{varwidth}
\usepackage{hyperref}
\usepackage{titlesec}
\usepackage{enumitem}
\usepackage{float}
\usepackage{caption}

\setlength{\parindent}{0pt}
\setlength{\parskip}{4pt}

\titleformat{\section}{\bfseries\Large}{}{0em}{}[\vspace{0.1cm}]
\titleformat{\subsection}{\bfseries\large}{}{0em}{}[\vspace{0.1cm}]

\hypersetup{
    colorlinks=true,
    linkcolor=black,
    urlcolor=blue,
    citecolor=black
}

\begin{document}

% ══════════════════════════════════════════════════════════════════════
%  PORTADA
% ══════════════════════════════════════════════════════════════════════
\begin{titlepage}
    \centering
    \vspace*{1.5cm}

    {\LARGE\bfseries Implementaci\'on de Autenticaci\'on\\[0.2cm]
     Multifactor (MFA) en Entornos Reales\par}

    \vspace{1.5cm}

    {\large\bfseries Jhoan Sebastian Parra\\[0.15cm]
     David Burbano Mari\~no}

    \vspace{1cm}

    {\large Seguridad Inform\'atica}

    \vspace{0.5cm}

    {\large Luz Marina Santos}

    \vspace{1.5cm}

    \includegraphics[width=4cm]{escudounipamplona.png}

    \vfill

    {\large Universidad de Pamplona\\[0.2cm]
     Junio 24 de 2026}

\end{titlepage}

% ══════════════════════════════════════════════════════════════════════
%  TABLA DE CONTENIDO
% ══════════════════════════════════════════════════════════════════════
\tableofcontents
\newpage

% ══════════════════════════════════════════════════════════════════════
%  1. TÍTULO Y DESCRIPCIÓN GENERAL
% ══════════════════════════════════════════════════════════════════════
\section{T\'itulo y Descripci\'on General}

\subsection{T\'itulo del Proyecto}

Implementaci\'on de Autenticaci\'on Multifactor (MFA) en el Simulador de Inversiones FIUP.

\subsection{Descripci\'on General}

El Simulador de Inversiones FIUP es una plataforma web educativa desarrollada para la Universidad de Pamplona. Permite a los estudiantes gestionar portafolios burs\'atiles virtuales, consultar cotizaciones en tiempo real mediante la API de Finnhub y participar en rankings acad\'emicos. La plataforma maneja datos personales (correo electr\'onico, n\'umero telef\'onico) y un saldo virtual de hasta 10,000 USD, lo que la convierte en un objetivo atractivo para ataques de suplantaci\'on de identidad.

El proyecto consisti\'o en integrar un sistema de Autenticaci\'on Multifactor (MFA) sobre la arquitectura existente del simulador. La implementaci\'on se realiz\'o sobre un backend en FastAPI con Python 3.12 y un frontend en React 18 con TypeScript. El MFA soporta tres modalidades de segundo factor: c\'odigos TOTP generados por aplicaciones como Google Authenticator, c\'odigos SMS enviados a trav\'es de una puerta de enlace HTTP, y c\'odigos de respaldo \'unicos pregenerados.

El flujo de inicio de sesi\'on se divide en dos pasos. Primero, el usuario ingresa su nombre de usuario y contrase\~na. Si las credenciales son v\'alidas y el usuario tiene habilitado el 2FA, el servidor emite un token temporal con validez de 30 segundos. En el segundo paso, el usuario debe presentar un c\'odigo TOTP de 6 d\'igitos, un c\'odigo SMS o uno de sus c\'odigos de respaldo. Solo entonces se emiten los tokens de acceso y actualizaci\'on definitivos.

% ══════════════════════════════════════════════════════════════════════
%  2. PROBLEMA Y JUSTIFICACIÓN
% ══════════════════════════════════════════════════════════════════════
\section{Problema Identificado y Justificaci\'on del Uso de MFA}

\subsection{Problema}

Las aplicaciones web financieras son el blanco principal de ataques de robo de credenciales. Un estudio del Verizon Data Breach Investigations Report de 2025 se\~nala que el 74\% de las brechas en aplicaciones web involucran credenciales d\'ebiles, robadas o reutilizadas. El Simulador FIUP, aunque es un entorno educativo, no est\'a exento de estos riesgos: almacena datos personales de estudiantes y gestiona activos virtuales que los usuarios valoran.

El problema central es que una contrase\~na, incluso con alta entrop\'ia, es un \'unico punto de fallo. Si un atacante obtiene la contrase\~na mediante phishing, keylogging, credential stuffing o una filtraci\'on de base de datos, obtiene acceso total a la cuenta de la v\'ictima. En el contexto del simulador, esto permite:
\begin{itemize}
    \item Manipular portafolios ajenos (compra y venta de activos).
    \item Robar datos personales de perfil (correo, tel\'efono).
    \item Suplantar la identidad del usuario en el ranking acad\'emico.
    \item Escalar a ataques contra otros usuarios mediante ingenier\'ia social.
\end{itemize}

\subsection{Justificaci\'on}

MFA mitiga este problema porque introduce un segundo factor que el atacante no posee. El modelo de seguridad cambia de ``algo que sabes'' (la contrase\~na) a ``algo que sabes + algo que tienes''. Incluso si la contrase\~na se ve comprometida, el atacante necesita acceso f\'isico al dispositivo del usuario para generar el c\'odigo TOTP, o haber interceptado el mensaje SMS correspondiente.

Seleccionamos TOTP como factor principal porque es un est\'andar abierto (RFC 6238), no requiere conectividad del usuario (el c\'odigo se genera localmente en el dispositivo) y no depende de operadores de telefon\'ia. Los c\'odigos SMS los incorporamos como factor de contingencia para usuarios sin acceso a un smartphone con aplicaci\'on autenticadora. Los c\'odigos de respaldo, almacenados como hash SHA-256, cubren el escenario de p\'erdida del dispositivo.

La implementaci\'on no se limita a agregar un segundo factor. Incluye medidas complementarias: rate limiting de 5 intentos por minuto en el endpoint de inicio de sesi\'on, tokens JWT con expiraci\'on corta (15 minutos para acceso, 7 d\'ias para actualizaci\'on), revocaci\'on de todos los tokens al cambiar la contrase\~na mediante un contador de versi\'on, y cookies HTTPOnly con SameSite=Lax para mitigar ataques XSS y CSRF.

% ══════════════════════════════════════════════════════════════════════
%  3. OBJETIVOS
% ══════════════════════════════════════════════════════════════════════
\section{Objetivo General y Objetivos Espec\'ificos}

\subsection{Objetivo General}

Implementar un sistema de autenticaci\'on multifactor funcional en el Simulador de Inversiones FIUP que combine contrase\~na y c\'odigos TOTP (RFC 6238) para reducir el riesgo de accesos no autorizados por robo de credenciales.

\subsection{Objetivos Espec\'ificos}

\begin{enumerate}[label=\textbf{O\arabic{enumi}}]
    \item Dise\~nar e implementar un flujo de registro con verificaci\'on de correo electr\'onico mediante c\'odigo OTP de 6 d\'igitos, almacenado en Redis con expiraci\'on de 10 minutos y un m\'aximo de 3 intentos de validaci\'on.
    \item Desarrollar un m\'odulo de generaci\'on y verificaci\'on de c\'odigos TOTP usando la librer\'ia pyotp, incluyendo la generaci\'on de c\'odigos QR para facilitar la configuraci\'on en aplicaciones autenticadoras.
    \item Implementar un sistema de tokens JWT con dos niveles: un token temporal de 30 segundos para el desaf\'io 2FA y tokens de acceso/actualizaci\'on con claims de seguridad (aud, iss, jti, type, \textit{password\_version}).
    \item Integrar rate limiting granular en todos los endpoints sensibles (login, registro, verificaci\'on 2FA, env\'io de SMS) usando slowapi con identificaci\'on por direcci\'on IP.
    \item Generar y almacenar 8 c\'odigos de respaldo por usuario, cifrados con SHA-256 y de un solo uso, para recuperaci\'on de cuenta en caso de p\'erdida del dispositivo.
    \item Construir una interfaz de usuario en React que gu\'ie al usuario a trav\'es de la configuraci\'on y el uso del 2FA, con manejo de estados de carga, error y expiraci\'on de token temporal.
\end{enumerate}

% ══════════════════════════════════════════════════════════════════════
%  4. DIAGRAMA DE ARQUITECTURA
% ══════════════════════════════════════════════════════════════════════
\section{Diagrama de Arquitectura Propuesta}

\subsection{Componentes del Sistema}

El sistema sigue una arquitectura de tres capas con un proxy inverso Nginx al frente. La siguiente tabla describe los componentes y su rol en el flujo MFA:

\begin{table}[H]
\centering
\caption{Componentes de la arquitectura MFA}
\begin{tblr}{colspec={l X[l]}, hlines, row{1}={font=\bfseries}}
\textbf{Componente} & \textbf{Rol en MFA} \\
Nginx (proxy) & Terminaci\'on SSL, rate limiting adicional, cabeceras de seguridad (HSTS, CSP, X-Frame-Options) \\
FastAPI (backend) & Orquestaci\'on del flujo MFA: validaci\'on de credenciales, emisi\'on de tokens, verificaci\'on TOTP, env\'io de OTP \\
PostgreSQL & Almacenamiento persistente: usuarios, secretos TOTP, c\'odigos de respaldo hasheados, c\'odigos de verificaci\'on \\
Redis & Almacenamiento temporal: datos de registro, tokens temporales de 2FA, contadores de rate limiting \\
Resend API & Entrega de c\'odigos OTP por correo electr\'onico durante el registro \\
SMS Gateway & Entrega de c\'odigos OTP por mensaje de texto durante el inicio de sesi\'on \\
React SPA (frontend) & Interfaz de usuario con manejo de estados 2FA, visualizaci\'on de c\'odigos QR, entrada de c\'odigos \\
\end{tblr}
\end{table}

\subsection{Flujo de Datos: Inicio de Sesi\'on con 2FA}

A continuaci\'on se describe el flujo paso a paso que debe reflejarse en el diagrama de arquitectura:

\begin{description}[itemsep=4pt]
    \item[Paso 1: Credenciales.] El usuario completa el formulario de inicio de sesi\'on con su nombre de usuario y contrase\~na. El frontend env\'ia una petici\'on POST a \texttt{/api/v1/login}.

    \item[Paso 2: Validaci\'on.] El backend consulta la tabla \textit{users} en PostgreSQL. Verifica la contrase\~na con bcrypt (12 rondas). Si el usuario tiene \textit{totp\_enabled = true}, no emite los tokens de acceso todav\'ia. En su lugar, genera un token JWT temporal con los siguientes atributos: tipo \textit{temp}, expiraci\'on de 30 segundos, y un identificador \'unico \textit{jti}. Almacena el \textit{jti} en Redis para garantizar que el token solo pueda usarse una vez. La respuesta incluye \textit{\{requires\_2fa: true, temp\_token, available\_methods\}}.

    \item[Paso 3: Selecci\'on de m\'etodo.] El frontend inspecciona \textit{available\_methods} y muestra las opciones disponibles (TOTP, SMS, c\'odigo de respaldo). Si solo hay un m\'etodo, avanza directamente.

    \item[Paso 4: Verificaci\'on TOTP.] El usuario ingresa el c\'odigo de 6 d\'igitos de su aplicaci\'on autenticadora. El frontend env\'ia una petici\'on POST a \texttt{/api/v1/2fa/login-verify} con \textit{\{temp\_token, code\}}.

    \item[Paso 5: Validaci\'on del segundo factor.] El backend decodifica el \textit{temp\_token} y verifica que el \textit{jti} exista en Redis y no haya sido usado. Luego lee el \textit{totp\_secret} del usuario desde PostgreSQL y ejecuta \textit{pyotp.TOTP(secret).verify(code, valid\_window=1)}. Si el c\'odigo es v\'alido, elimina el \textit{jti} de Redis (consumo del token temporal). Genera un \textit{access\_token} con expiraci\'on de 15 minutos y un \textit{refresh\_token} con expiraci\'on de 7 d\'ias. Ambos incluyen el \textit{password\_version} actual del usuario.

    \item[Paso 6: Entrega de sesi\'on.] El backend establece cookies HTTPOnly con los tokens y devuelve el objeto de usuario en el cuerpo de la respuesta. El frontend almacena el estado de autenticaci\'on en Zustand (solo el objeto de usuario, nunca los tokens) y redirige al tablero principal.
\end{description}

\subsection{Flujo de Datos: Configuraci\'on de 2FA}

\begin{description}[itemsep=4pt]
    \item[Paso 1: Solicitud de configuraci\'on.] El usuario autenticado accede a la p\'agina de configuraci\'on y hace clic en ``Habilitar 2FA''. El frontend llama a \texttt{GET /api/v1/2fa/setup}.

    \item[Paso 2: Generaci\'on de secreto.] El backend genera un secreto TOTP de 16 caracteres en base32 usando \textit{pyotp.random\_base32()}. Construye la URI de aprovisionamiento seg\'un el formato \texttt{otpauth://totp/...}. Genera un c\'odigo QR en formato PNG y lo codifica en base64. Devuelve \textit{\{secret, qr\_code, provisioning\_uri\}}. El secreto a\'un no se persiste en la base de datos.

    \item[Paso 3: Escaneo y verificaci\'on.] El usuario escanea el c\'odigo QR con su aplicaci\'on autenticadora (Google Authenticator, Authy, etc.) e ingresa el c\'odigo de 6 d\'igitos que la aplicaci\'on genera. El frontend env\'ia \texttt{POST /api/v1/2fa/verify \{code\}}.

    \item[Paso 4: Activaci\'on.] El backend verifica el c\'odigo. Si es v\'alido, persiste el secreto en la columna \textit{totp\_secret} del usuario en PostgreSQL, establece \textit{totp\_enabled = true} y genera 8 c\'odigos de respaldo. Cada c\'odigo de respaldo tiene el formato \textit{XXXXXX-XXXXXX} y se almacena como hash SHA-256 en la tabla \textit{backup\_codes}. La respuesta incluye los c\'odigos de respaldo en texto plano para que el usuario los guarde.
\end{description}

% ══════════════════════════════════════════════════════════════════════
%  5. HERRAMIENTAS
% ══════════════════════════════════════════════════════════════════════
\section{Herramientas Tecnol\'ogicas Seleccionadas}

La siguiente tabla resume las herramientas utilizadas en la implementaci\'on del MFA, junto con su prop\'osito espec\'ifico dentro del flujo de autenticaci\'on.

\begin{table}[H]
\centering
\caption{Herramientas tecnol\'ogicas del proyecto MFA}
\begin{tblr}{colspec={p{3.5cm} p{1.8cm} X[l]}, hlines, row{1}={font=\bfseries}}
\textbf{Herramienta} & \textbf{Versi\'on} & \textbf{Prop\'osito en MFA} \\
Python 3.12 + FastAPI & 0.135 & Backend as\'incrono del flujo MFA. Maneja solicitudes concurrentes sin bloquear el event loop durante verificaci\'on de OTP o consultas a Redis \\
bcrypt & 5.0+ & Hashing de contrase\~nas con 12 rondas. Resistente a ataques con GPU y ASIC. Supera el m\'inimo recomendado por NIST de 10 rondas \\
PyJWT & 2.12 & Creaci\'on y verificaci\'on de tokens JWT. Soporta claims personalizados (aud, iss, jti, type) para prevenir re\'usos entre servicios \\
pyotp & 2.9 & Implementaci\'on del est\'andar TOTP RFC~6238. Generaci\'on de secretos, c\'odigos QR y verificaci\'on con ventana de tolerancia \\
qrcode[pil] & 7.4 & Generaci\'on de c\'odigos QR en servidor para facilitar la configuraci\'on del 2FA sin transferir secretos por canales inseguros \\
slowapi & 0.1.9 & Rate limiting por direcci\'on IP. L\'imites diferenciados: 5/min en login, 3/min en env\'io de SMS, 10/min en verificaci\'on 2FA \\
Redis & 7.0 & Almacenamiento en memoria de datos ef\'imeros: tokens temporales de 2FA, datos de registro, contadores de rate limiting. TTL autom\'atico \\
Resend & 0.5 & API de env\'io de correos electr\'onicos para c\'odigos OTP de verificaci\'on de registro \\
httpx & 0.28 & Cliente HTTP as\'incrono para la comunicaci\'on con la puerta de enlace SMS externa \\
React 18 + TypeScript & 18/5.3 & Interfaz de usuario con manejo de estados 2FA, visualizaci\'on de c\'odigos QR y entrada de c\'odigos con validaci\'on Zod \\
Zustand & 4 & Gesti\'on de estado del lado del cliente. Persiste solo el objeto de usuario en localStorage, nunca los tokens \\
Zod & \textendash & Validaci\'on de esquemas en el frontend: formato de c\'odigos, complejidad de contrase\~nas, campos obligatorios \\
\end{tblr}
\end{table}

\subsection{Criterios de Selecci\'on}

Elegimos estas herramientas por cuatro razones principales. Primera, la compatibilidad con el stack existente del simulador (FastAPI + React), lo que evit\'o reescribir la infraestructura. Segunda, la madurez de las bibliotecas: pyotp tiene m\'as de 10 a\~nos de desarrollo continuo y es la implementaci\'on Python de referencia para TOTP. Tercera, el rendimiento as\'incrono: tanto FastAPI como httpx y redis-py operan sobre el mismo event loop, eliminando cuellos de botella en la verificaci\'on de c\'odigos. Cuarta, el costo: todas las herramientas son de c\'odigo abierto y no requieren licencias adicionales.

% ══════════════════════════════════════════════════════════════════════
%  6. EVALUACIÓN DE RESULTADOS
% ══════════════════════════════════════════════════════════════════════
\section{Evaluaci\'on de Resultados}

\subsection{C\'omo MFA Mejora la Seguridad}

La tabla siguiente compara el estado de seguridad del sistema antes y despu\'es de la implementaci\'on del MFA:

\begin{table}[H]
\centering
\caption{Evaluaci\'on de mejora de seguridad con MFA}
\begin{tblr}{colspec={p{3.5cm} p{2.5cm} p{3.2cm} X[l]}, hlines, row{1}={font=\bfseries}}
\textbf{M\'etrica} & \textbf{Sin MFA} & \textbf{Con MFA} & \textbf{Impacto} \\
Factores de autenticaci\'on & 1 (contrase\~na) & 2 (contrase\~na + TOTP/SMS/respaldo) & El atacante necesita ambos factores \\
Contrase\~na comprometida $\rightarrow$ acceso total & S\'i & No & Reduce ventana de ataque \\
Fuerza bruta en login & Sin l\'imite efectivo & 5 intentos/min por IP + segundo paso & Requiere m\'ultiples IPs para mantener 1 intento/s \\
Token interceptado v\'alido por & 30 minutos & 15 min + \textit{password\_version} & Ventana de ataque reducida 50\% \\
Registro de cuentas falsas & Sin verificaci\'on & OTP obligatorio por email & Previene creaci\'on masiva automatizada \\
Sesi\'on secuestrada por XSS & Posible (token en localStorage) & Cookies HTTPOnly + SameSite=Lax & Token inaccesible desde JavaScript \\
\end{tblr}
\end{table}

El mayor impacto se da en el escenario de credential stuffing. Un atacante que obtiene una lista de contrase\~nas filtradas no puede iniciar sesi\'on sin tener adem\'as el dispositivo del usuario. Incluso si logra evadir el rate limit usando una botnet, debe resolver el desaf\'io TOTP, que expira en 30 segundos y es de un solo uso.

\subsection{Limitaciones de la Implementaci\'on Actual}

Ning\'un sistema de seguridad es perfecto. Identificamos varias limitaciones que deben considerarse:

\begin{itemize}
    \item \textbf{Secreto TOTP en texto plano.} El campo \textit{totp\_secret} en la base de datos PostgreSQL no est\'a cifrado. Si un atacante obtiene un volcado de la base de datos, mediante inyecci\'on SQL, acceso f\'isico o una copia de seguridad mal protegida, puede generar c\'odigos TOTP v\'alidos sin l\'imite para cualquier usuario. La soluci\'on pasa por cifrar el secreto con Fernet (cryptography sim\'etrica) usando una clave derivada del \textit{SECRET\_KEY} de la aplicaci\'on.
    \item \textbf{El cierre de sesi\'on no invalida los tokens.} La funci\'on de logout actualmente solo elimina las cookies del navegador. Los tokens JWT subyacentes siguen siendo v\'alidos hasta su expiraci\'on natural (15 minutos para access token, 7 d\'ias para refresh token). Un atacante que haya interceptado un token puede seguir us\'andolo incluso despu\'es de que el usuario cierre la sesi\'on. La soluci\'on es mantener una lista negra de tokens en Redis (\textit{blacklist:\{jti\}}) y verificar contra ella en cada solicitud.
    \item \textbf{Ausencia de CAPTCHA.} El rate limiting por IP es efectivo contra atacantes individuales, pero una botnet distribuida puede eludirlo f\'acilmente usando miles de direcciones IP diferentes. Agregar Google reCAPTCHA v3 o hCaptcha en los formularios de login y registro a\~nadir\'ia una capa de defensa basada en an\'alisis de comportamiento.
    \item \textbf{Dependencia de SMS como factor.} La NIST ya no recomienda SMS como segundo factor desde la publicaci\'on del SP 800-63 en 2017. Los c\'odigos SMS son vulnerables a ataques de SIM swapping y a la interceptaci\'on de la se\~nal SS7. En nuestro dise\~no lo mantenemos como factor de contingencia, no como factor principal, pero la recomendaci\'on es migrar gradualmente a WebAuthn (FIDO2/Passkeys).
    \item \textbf{Sin registro de auditor\'ia.} El sistema no genera registros de eventos 2FA: qui\'en habilit\'o la autenticaci\'on, desde qu\'e direcci\'on IP, cu\'ando se us\'o un c\'odigo de respaldo, o cu\'antos intentos fallidos de 2FA han ocurrido. Esta informaci\'on es necesaria para el an\'alisis forense en caso de incidente.
\end{itemize}

\subsection{Propuestas de Mejora Futura}

Las limitaciones anteriores se traducen en un plan de mejoras organizado en tres horizontes temporales:

\textbf{Corto plazo (pr\'oximos 30 d\'ias):}
\begin{itemize}
    \item Cifrar el secreto TOTP con la biblioteca \textit{cryptography.fernet}. La clave de cifrado se deriva del \textit{SECRET\_KEY} de la aplicaci\'on usando HKDF.
    \item Implementar una lista negra de tokens JWT en Redis. Al cerrar sesi\'on, el \textit{jti} del token se agrega a la lista negra con un TTL igual al tiempo restante de expiraci\'on del token. Cada solicitud entrante verifica que su \textit{jti} no est\'e en la lista negra.
    \item Agregar logging estructurado para eventos de 2FA: habilitaci\'on, deshabilitaci\'on, inicio de sesi\'on exitoso y uso de c\'odigo de respaldo.
\end{itemize}

\textbf{Mediano plazo (1 a 3 meses):}
\begin{itemize}
    \item Integrar Google reCAPTCHA v3 en los formularios de inicio de sesi\'on y registro. reCAPTCHA v3 funciona en segundo plano sin interrumpir al usuario, asignando una puntuaci\'on de riesgo que el backend puede usar para denegar solicitudes sospechosas.
    \item Implementar WebAuthn (FIDO2/Passkeys) como factor adicional. A diferencia del TOTP, WebAuthn es resistente a ataques de phishing porque vincula la credencial al origen del sitio web.
    \item Agregar un panel de actividad de cuenta donde el usuario pueda ver los \'ultimos inicios de sesi\'on con direcci\'on IP, dispositivo y m\'etodo 2FA utilizado.
\end{itemize}

\textbf{Largo plazo (pr\'oximo ciclo acad\'emico):}
\begin{itemize}
    \item Deprecar el SMS como factor 2FA y reemplazarlo completamente por notificaciones push a una aplicaci\'on m\'ovil dedicada.
    \item Implementar un sistema de detecci\'on de anomal\'ias que analice patrones de inicio de sesi\'on (ubicaci\'on geogr\'afica, horario, agente de usuario) y solicite autenticaci\'on adicional cuando el riesgo sea alto.
\end{itemize}

\subsection{Desaf\'ios de Ingenier\'ia Encontrados}

Durante el desarrollo enfrentamos problemas que vale la pena documentar. El manejo de la expiraci\'on del token temporal fue el m\'as complejo: si el usuario se demora m\'as de 30 segundos en ingresar el c\'odigo TOTP, el token expira y debe comenzar de nuevo. En la interfaz de usuario, esto se traduce en un mensaje de ``Sesi\'on expirada'' y un redireccionamiento al paso de credenciales. Al principio no manej\'abamos este caso y el usuario ve\'a un error gen\'erico.

Otro desaf\'io fue la ventana de tolerancia de reloj. El est\'andar TOTP asume que el reloj del dispositivo del usuario est\'a sincronizado con el servidor. En pruebas encontramos que algunos tel\'efonos ten\'ian una diferencia de hasta 45 segundos, lo que provocaba falsos rechazos. Solucionamos esto configurando \textit{valid\_window=1} en pyotp, que acepta el c\'odigo del intervalo anterior y el siguiente adem\'as del actual. Esto ampl\'ia la ventana a 90 segundos en total, un compromiso aceptable entre seguridad y usabilidad.

La gesti\'on de sesiones comprometidas requiri\'o el mecanismo de \textit{password\_version}. Cada vez que el usuario cambia su contrase\~na, este contador se incrementa y todos los tokens emitidos con la versi\'on anterior se invalidan autom\'aticamente. Implementamos esto en \textit{auth\_service.py} y lo verificamos en cada solicitud protegida.

% ══════════════════════════════════════════════════════════════════════
%  7. CONCLUSIONES
% ══════════════════════════════════════════════════════════════════════
\section{Conclusiones Generales}

La implementaci\'on de MFA en el Simulador de Inversiones FIUP demuestra que es posible integrar autenticaci\'on multifactor en una aplicaci\'on web existente sin reescribir la arquitectura base. El sistema soporta tres modalidades de segundo factor (TOTP, SMS y c\'odigos de respaldo), lo que cubre tanto el escenario principal como las situaciones de contingencia.

Desde el punto de vista t\'ecnico, la implementaci\'on cumple con los est\'andares de seguridad establecidos: bcrypt con 12 rondas para contrase\~nas, JWT con validaci\'on de audiencia y emisor, tokens temporales de un solo uso, rate limiting granular y revocaci\'on de sesi\'on por cambio de contrase\~na. El uso de Redis como almacenamiento de datos ef\'imeros permiti\'o manejar los tokens temporales y los c\'odigos OTP con tiempos de expiraci\'on precisos sin contaminar la base de datos relacional.

La evaluaci\'on de resultados muestra una mejora significativa en la postura de seguridad. Un atacante que comprometa la contrase\~na de un usuario todav\'ia necesita acceso al dispositivo f\'isico para completar la autenticaci\'on. El rate limiting reduce la tasa de ataques de fuerza bruta a 5 intentos por minuto por direcci\'on IP. La combinaci\'on de cookies HTTPOnly, SameSite=Lax y la ausencia de tokens en localStorage elimina vectores de ataque XSS y CSRF.

Sin embargo, el sistema tiene limitaciones que no deben ignorarse. La m\'as cr\'itica es el almacenamiento del secreto TOTP en texto plano en la base de datos. Le sigue la falta de invalidaci\'on de tokens al cerrar sesi\'on. Ambas son corregibles en el corto plazo y no invalidan el prototipo funcional. La ausencia de CAPTCHA y de registro de auditor\'ia son deficiencias que deben abordarse antes de un despliegue en producci\'on.

Como trabajo futuro, la migraci\'on a WebAuthn (FIDO2/Passkeys) representa el siguiente paso l\'ogico. Este est\'andar elimina la dependencia de c\'odigos de 6 d\'igitos y ofrece resistencia al phishing, algo que el TOTP no puede garantizar. La integraci\'on de un sistema de detecci\'on de anomal\'ias basado en patrones de inicio de sesi\'on permitir\'ia escalar din\'amicamente los requisitos de autenticaci\'on seg\'un el nivel de riesgo.

En conclusi\'on, el prototipo de MFA implementado es funcional, sigue buenas pr\'acticas de seguridad y cumple con los objetivos acad\'emicos planteados. Con las correcciones se\~naladas, constituye una base s\'olida para un sistema de autenticaci\'on robusto en un entorno educativo y, con mejoras adicionales, para un despliegue en producci\'on.

\end{document}
